\section*{Capítulo \ref{ch:b3:atmospheric-dynamics} --- Dinâmica atmosférica}

\begin{solution}{exo:b3:atmospheric-dynamics:1}
(a)~$\zeta = 2\omega = 2\times2\pi/86\,400 = \qty{1.45e-4}{s^{-1}}$.
(b)~$f = \num{7.29e-5}$ e \qty{1.26e-4}{s^{-1}}: o vórtice tem cerca do
dobro da \omterm{def:b3:atmospheric-dynamics:vorticity-kinds}{vorticidade planetária} a \ang{30} N e um pouco mais que a de
\ang{60} N.
\end{solution}

\begin{solution}{exo:b3:atmospheric-dynamics:2}
$\zeta = f_{30} - f_{40} = \num{7.29e-5} - \num{9.37e-5} =
\qty{-2.08e-5}{s^{-1}}$: anticiclônico.
\end{solution}

\begin{solution}{exo:b3:atmospheric-dynamics:3}
$k = 2\pi/(5\times10^6) = \qty{1.257e-6}{m^{-1}}$; $c = -\beta/k^2 =
-\num{1.62e-11}/\num{1.58e-12} = \qty{-10.3}{m/s}$: para oeste.
\end{solution}

\begin{solution}{exo:b3:atmospheric-dynamics:4}
$\beta = 2\Omega\cos\ang{60}/a = \qty{1.14e-11}{m^{-1}.s^{-1}}$;
$\lambda_s = 2\pi\sqrt{10/\num{1.14e-11}} = \qty{5870}{km}$; o círculo,
\qty{20015}{km}, comporta $n_s = 3.4$ ondas.
\end{solution}

\begin{solution}{exo:b3:atmospheric-dynamics:5}
$q = f/h$ é conservado: sobre a serra $\zeta = f(8/10) - f =
\qty{-2.06e-5}{s^{-1}}$, anticiclônico; sobre a bacia $f(12/10) - f =
\qty{+2.06e-5}{s^{-1}}$, ciclônico. Onde as colunas são esticadas, a sotavento
da serra e sobre a bacia, cria-se vorticidade ciclônica: ali
tendem a se formar depressões.
\end{solution}

\begin{solution}{exo:b3:atmospheric-dynamics:6}
$k = 5/(a\cos\ang{45}) = \qty{1.11e-6}{m^{-1}}$, $\beta/k^2 =
\qty{13.1}{m/s}$: $c = \qty{1.9}{m/s}$, $c_g = \qty{28.1}{m/s}$. A energia
percorre \qty{5000}{km} em \num{2.1}~dias; o padrão, em \num{31}~dias.
\end{solution}

\begin{solution}{exo:b3:atmospheric-dynamics:7}
(a)~$k^2 = \beta/U - l^2 = \num{1.079e-12} - \num{2.74e-13} =
\qty{8.05e-13}{m^{-2}}$, $k = \qty{8.97e-7}{m^{-1}}$: $\lambda =
\qty{7000}{km}$ e $n = 4.0$, menos ondas do que com $l = 0$. (b)~$l^2 =
\num{1.097e-12} > \beta/U = \num{1.079e-12}$: $k^2$ seria negativo.
\end{solution}

\begin{solution}{exo:b3:atmospheric-dynamics:8}
$L_d = 0.012\times10^4/\num{1.031e-4} = \qty{1164}{km}$; $\sigma =
0.31\times\num{1.031e-4}\times\num{2e-3}/0.012 = \qty{5.3e-6}{s^{-1}}$, um
tempo de $e$-dobramento de \num{2.2}~dias; comprimento de onda $2\pi\times1164/1.61 =
\qty{4540}{km}$.
\end{solution}

\begin{solution}{exo:b3:atmospheric-dynamics:9}
$[v^*T^*] = (VT_0/2)\cos\delta$: \qty{25}{K.m/s}, \qty{12.5}{K.m/s} e 0;
fluxos de calor \num{25100}, \num{12600} e \qty{0}{W/m^2}. Quando a
temperatura está defasada de um quarto de comprimento de onda em relação ao
vento, tanto ar quente vai para o sul quanto para o norte: não há transporte.
\end{solution}

\begin{solution}{exo:b3:atmospheric-dynamics:10}
Linearizando em torno de $U$, com $\zeta' = \nabla^2\psi'$ e $v' =
\partial_x\psi'$: $(\partial_t + U\partial_x)\nabla^2\psi' + \beta\partial_x
\psi' = 0$. Com $\psi' \propto \mathrm e^{\mathrm i(kx + ly - \omega t)}$,
$(\omega - Uk)K^2 = -\beta k$, logo $\omega = Uk - \beta k/K^2$. Em relação
ao ar, a velocidade de fase é $c - U = -\beta/K^2 < 0$ para toda onda.
\end{solution}

\begin{solution}{exo:b3:atmospheric-dynamics:11}
$\varphi_H = [5\times9.81\times12\,000\times0.25/(3\times(\num{7.292e-5})^2
\times(\num{6.371e6})^2)]^{1/2} = \qty{0.477}{rad}$, \ang{27.3}.
$u_M = \Omega a\varphi_H^2 = \qty{106}{m/s}$; exatamente, $\Omega a\sin^2
\varphi_H/\cos\varphi_H = \qty{110}{m/s}$: duas vezes e meia a
corrente de jato observada.
\end{solution}

\begin{solution}{exo:b3:atmospheric-dynamics:12}
(a)~Ao longo da faixa, a interface é deslocada de $\eta = sy$, $-L/2 \le y
\le L/2$, a partir de sua posição plana. Uma camada de contraste de
densidade reduzido, deslocada de $\eta$ armazena $\tfrac12\rho g'\eta^2$ por unidade de área; média
sobre a faixa, $\langle\eta^2\rangle = s^2L^2/12$. (b)~$g' = 0.327$:
\num{1.5e5} e \qty{6.1e5}{J/m^2}, contra $\tfrac12\times0.7\times15^2
\times10^4 = \qty{7.9e5}{J/m^2}$ de energia cinética: comparáveis.
\end{solution}

\begin{solution}{pb:b3:atmospheric-dynamics:1}
\textbf{1.} $f_{45} = \qty{1.031e-4}{s^{-1}}$, $f_{50} =
\qty{1.117e-4}{s^{-1}}$; $\zeta = f_{45} - f_{50} = \qty{-8.6e-6}{s^{-1}}$.

\textbf{2.} $R = V/|\zeta| = 15/\num{8.6e-6} = \qty{1750}{km}$; no sentido horário,
de volta para o sul.

\textbf{3.} $\zeta = f_{45}(9.5/10) - f_{50} = \qty{-1.38e-5}{s^{-1}}$, mais
anticiclônico.

\textbf{4.} $\zeta = f_{45} - f_{40} = \qty{+9.4e-6}{s^{-1}}$, ciclônico: ela
volta para o norte. Deslocada em qualquer dos sentidos, a parcela adquire a
vorticidade que a faz retornar: um efeito restaurador, e uma oscilação.

\textbf{5.} $\beta = \qty{1.62e-11}{m^{-1}.s^{-1}}$; o círculo tem
\qty{28306}{km} de comprimento.

\textbf{6.} $n = 3$ a 8: comprimentos de onda 9435, 7076, 5661, 4718, 4044 e
\qty{3538}{km}; $c = -21.5$, $-5.5$, 1.9, 5.9, 8.3 e \qty{9.9}{m/s}; $c_g
= 51.5$, 35.5, 28.1, 24.1, 21.7 e \qty{20.1}{m/s}.

\textbf{7.} $n = 3$ e 4 movem-se para oeste, $n \ge 5$ para leste. A onda $n = 5$ leva
$10^6/1.86 = \qty{5.4e5}{s}$, \num{6.2}~dias, para percorrer \qty{1000}{km}.

\textbf{8.} A jusante, para leste: com $c_g = \qty{24.1}{m/s}$ a energia percorre
\qty{5000}{km} em \num{2.4}~dias.

\textbf{9.} $\lambda_s = 2\pi\sqrt{15/\num{1.62e-11}} = \qty{6050}{km}$,
$n_s = 4.7$.

\textbf{10.} $n_s = 5.7$ para \qty{10}{m/s} e 4.1 para \qty{20}{m/s}: um
vento de oeste mais forte favorece ondas estacionárias mais longas, em menor
número ao redor do hemisfério.

\textbf{11.} $n_s = 5.8$ em \ang{35} N e 3.4 em \ang{55} N.

\textbf{12.} $\zeta = f(8/10) - f = \qty{-2.06e-5}{s^{-1}}$: o estreitamento
sobre as montanhas torna a coluna anticiclônica; o alongamento a sotavento,
ciclônico. As montanhas impõem um padrão fixo de vorticidade ao escoamento,
um forçamento que não se desloca.

\textbf{13.} Um forçamento fixo produz uma resposta permanente; a resposta é
máxima quando o padrão forçado é ele próprio uma onda livre que ficaria
parada, a ressonância em $\lambda_s$.

\textbf{14.} $\Lambda = 25/10^4 = \qty{2.5e-3}{s^{-1}}$;
$|\partial T/\partial y| = \Lambda fT/g = \num{2.5e-3}\times\num{1.031e-4}
\times250/9.81 = \qty{6.6}{K}$ por \qty{1000}{km}.

\textbf{15.} $\mathrm{Ri} = 0.011^2/(\num{2.5e-3})^2 = 19$: muito acima de 0.25,
sem turbulência de cisalhamento.

\textbf{16.} $L_d = 0.011\times10^4/\num{1.031e-4} = \qty{1067}{km}$.

\textbf{17.} $2\pi\times1067/1.61 = \qty{4160}{km}$; $28306/4160 = 6.8$
ondas.

\textbf{18.} $2\pi\times1067/2.40 = \qty{2790}{km}$.

\textbf{19.} $\sigma = 0.31\times\num{1.031e-4}\times\num{2.5e-3}/0.011 =
\qty{7.3e-6}{s^{-1}}$, \qty{0.63}{dia^{-1}}.

\textbf{20.} $\exp(3\times0.627) = 6.6$.

\textbf{21.} A tempestade, com \qty{4160}{km}, é menor do que o comprimento de
onda estacionário, \qty{6050}{km}: como onda de Rossby, ela é levada para leste
pelos ventos de oeste, e as tempestades se deslocam para leste pelas ondas
longas mais lentas ou estacionárias.

\textbf{22.} Inclinação para oeste com a altura: só então o ar que se move para os polos é
mais quente que o ar que se move para o equador, de modo que o calor segue para os polos e a
\omterm{def:b3:atmospheric-dynamics:ape}{energia potencial disponível} é liberada.

\textbf{23.} A perturbação de crescimento mais rápido cresce por um fator $e$ a cada
\textbf{\num[text-series-to-math=true]{1.6}~dias} (\num{1.59}~dias pelo
cálculo).
\end{solution}
